\subsection{CONSTRUCTION OF CAYLEY ALGEBRAS. QUATERNIONS}

Let \((\mathbf{E}, s)\) be a Cayley algebra over A, for which we shall use the notation of no. 4, and let \(\gamma \in \mathbf{A}\) . Let F be the algebra over A whose underlying module is \(\mathbf{E} \times \mathbf{E}\) and whose multiplication is defined by

\[
(x, y) (x ^ {\prime}, y ^ {\prime}) = (x x ^ {\prime} + \gamma \bar {y} ^ {\prime} y, y \bar {x} ^ {\prime} + y ^ {\prime} x);\tag{22}
\]

clearly \((e,0)\) is unit element of F and \(E \times \{0\}\) is a subalgebra of F isomorphic to E; we shall identify it with E in what follows, so that \(x \in E\) is identified with \((x,0)\) and in particular e is identified with the unit element of F.

Let \(t\) be the permutation of \(\mathbf{F}\) defined by

\begin{enumerate}
\def\labelenumi{(\arabic{enumi})}
\setcounter{enumi}{22}
\tightlist
\item
\end{enumerate}

\[
t ((x, y)) = (\bar {x}, - y) \quad (x \in \mathbf {E}, y \in \mathbf {E}).
\]

PROPOSITION 5. (i) The ordered pair \((\mathbf{F}, t)\) is a Cayley algebra over \(\mathbf{A}\) .

\begin{enumerate}
\def\labelenumi{(\roman{enumi})}
\setcounter{enumi}{1}
\tightlist
\item
  Let \(j = (0, e)\) , so that \((x, y) = xe + yj\) for \(x \in \mathbf{E}\) , \(y \in \mathbf{E}\) . The Cayley trace and norm \(\mathbf{T}_{\mathbb{F}}\) and \(\mathbf{N}_{\mathbb{F}}\) of \(\mathbf{F}\) are given by the formulae
\end{enumerate}

\[
\mathrm{T} _ {\mathrm{F}} (x e + y j) = \mathrm{T} (x), \quad \mathrm{N} _ {\mathrm{F}} (x e + y j) = \mathrm{N} (x) - \gamma \mathrm{N} (y).\tag{24}
\]

\begin{enumerate}
\def\labelenumi{(\roman{enumi})}
\setcounter{enumi}{2}
\tightlist
\item
  For F to be associative, it is necessary and sufficient that E be associative and commutative.
\end{enumerate}

For \((x,y)\in \mathbf{F},\)

\begin{enumerate}
\def\labelenumi{(\arabic{enumi})}
\setcounter{enumi}{24}
\tightlist
\item
\end{enumerate}

\[
\begin{array}{r l r} & & {(x, y) + t ((x, y)) = (x + \bar {x}, 0) = \mathrm{T} (x) e} \\ & & {(x, y) t ((x, y)) = (x, y) (\bar {x}, - y) = (x \bar {x} - \gamma \bar {y} y, y \bar {\bar {x}} - y x)} \\ & & {= (\mathrm{N} (x) - \gamma \mathrm{N} (y), 0) = (\mathrm{N} (x) - \gamma \mathrm{N} (y)) e.} \end{array}\tag{26}
\]

To prove both (i) and (ii), it therefore suffices to show that t is an antiautomorphism of F. Clearly t is an A-linear bijection. On the other hand,

\[
\begin{array}{r l} t ((x, y) \cdot (x ^ {\prime}, y ^ {\prime})) & = t ((x x ^ {\prime} + \gamma \bar {y} ^ {\prime} y, y \bar {x} ^ {\prime} + y ^ {\prime} x)) = (\bar {x} ^ {\prime} \bar {x} + \gamma \bar {y} y ^ {\prime}, - y \bar {x} - y ^ {\prime} x) \\ & = (\bar {x} ^ {\prime}, - y ^ {\prime}) (\bar {x}, - y) = t ((x ^ {\prime}, y ^ {\prime})) t ((x, y)) \end{array}
\]

and hence \(t\) is an antiautomorphism.

It remains to prove (iii). As E is identified with a subalgebra of F, E may be assumed to be associative. Let \(u = (x, y)\) , \(u' = (x', y')\) , \(u'' = (x'', y'')\) be elements of F. Then

\[
\left\{ \begin{array}{l} (u u ^ {\prime}) u ^ {\prime \prime} = ((x x ^ {\prime} + \gamma \bar {y} ^ {\prime} y) x ^ {\prime \prime} + \gamma \bar {y} ^ {\prime \prime} (\bar {y} x ^ {\prime} + y ^ {\prime} x), \\ \qquad \qquad \qquad \qquad \qquad \qquad \qquad \qquad \qquad \qquad (y \bar {x} ^ {\prime} + y ^ {\prime} x) \bar {x} ^ {\prime \prime} + y ^ {\prime \prime} (x x ^ {\prime} + \gamma \bar {y} ^ {\prime} y)) \\ u (u ^ {\prime} u ^ {\prime \prime}) = (x (x ^ {\prime} x ^ {\prime \prime} + \gamma \bar {y} ^ {\prime \prime} y ^ {\prime}) + (\gamma (x ^ {\prime \prime} \bar {y} ^ {\prime} + \bar {x} ^ {\prime} \bar {y} ^ {\prime \prime}) y, \\ \qquad \qquad \qquad \qquad \qquad \qquad \qquad \qquad \qquad \qquad y (\bar {x} ^ {\prime \prime} \bar {x} ^ {\prime} + \gamma \bar {y} ^ {\prime} y ^ {\prime \prime}) + (y ^ {\prime} \bar {x} ^ {\prime \prime} + y ^ {\prime \prime} x ^ {\prime}) x). \end{array} \right.\tag{27}
\]

Examining these formulae shows that the commutativity of E implies the associativity of F. Conversely, if F is associative, formulae (27) applied with \(y = y' = 0\) , \(x'' = 0\) and \(y'' = e\) give \((0, x'x) = (0, xx')\) , that is \(x'x = xx'\) for all \(x, x'\) in E; thus E is then commutative.

Note also that, in the above notation, for \(x, y\) in \(\mathbf{E}\) ,

\begin{enumerate}
\def\labelenumi{(\arabic{enumi})}
\setcounter{enumi}{27}
\tightlist
\item
\end{enumerate}

\[
\begin{array}{r l} y j = j \overline {{y}}, & x (y j) = (y x) j, \quad (x j) y = (x \overline {{y}}) j, \quad (x j) (y j) = \overline {{y}} x e \\ & j ^ {2} = e. \end{array}\tag{29}
\]

The Cayley algebra \((\mathbf{F}, t)\) is called the Cayley extension of \((\mathbf{E}, s)\) defined by \(\gamma\) .

Examples. (1) If we take \(\mathbf{E} = \mathbf{A}\) (and hence \(s = 1_{\mathbf{A}}\) ), the algebra \(\mathbf{F}\) is a quadratic \(\mathbf{A}\) -algebra with basis \((e,j)\) where \(j^2 = \gamma e\) .

\begin{enumerate}
\def\labelenumi{(\arabic{enumi})}
\setcounter{enumi}{1}
\tightlist
\item
  Take E to be a quadratic algebra of type \((\alpha, \beta)\) , so that the underlying module of E is \(A^{2}\) , with multiplication table (2) (no. 3) for the canonical basis. Take s to be conjugation of E (no. 3, Proposition 2). Then, for all \(\gamma \in A\) , the Cayley extension F of \((E, s)\) defined by \(\gamma\) is called the quaternion algebra of type \((\alpha, \beta, \gamma)\) , which is associative by no. 3, Proposition 1 and Proposition 5 above; its underlying module is \(A^{4}\) and, if \((e, i, j, k)\) denotes the canonical basis of \(A^{4}\) , the corresponding multiplication table is given by
\end{enumerate}

\[
\left\{ \begin{array}{l l} i ^ {2} = \alpha e + \beta i, & i j = k, \qquad i k = \alpha j + \beta k, \\ j i = \beta j - k, & j ^ {2} = \gamma e, \qquad j k = \beta \gamma e - \gamma i, \\ k i = - \alpha j, & k j = \gamma i, \qquad k ^ {2} = - \alpha \gamma e. \end{array} \right.\tag{30}
\]

Further, for \(u = \rho e + \xi i + \eta j + \eta k\) (with \(\rho, \xi, \eta, \zeta\) in A), we have (writing \(\bar{u}\) instead of \(t(u)\) and identifying A with Ae):

\[
\left\{ \begin{array}{c} u = (\rho + \beta \xi) e - \xi i - \eta j - \zeta k \\ \mathrm{T} _ {\mathrm{F}} (u) = 2 \rho + \beta \xi \\ \mathrm{N} _ {\mathrm{F}} (u) = \rho^ {2} + \beta \rho \xi - \alpha \xi^ {2} - \gamma (\eta^ {2} + \beta \eta \zeta - \alpha \zeta^ {2}). \end{array} \right.\tag{31}
\]

Formulae (30) follow from (28) and (29) and formulae (31) from (23) and (24), taking account of the formulae for the quadratic algebra E.

Then, for \(u, v\) in \(\mathbf{F}\) ,

\[
\mathrm{N} _ {\mathrm{F}} (u v) = \mathrm{N} _ {\mathrm{F}} (u) \mathrm{N} _ {\mathrm{F}} (v)\tag{32}
\]

for \(\mathrm{N}_{\mathbb{F}}(uv)=uv.\overline{uv}=uv(\bar{v}.\bar{u})=u(v\bar{v})\bar{u}=(u\bar{u})(v\bar{v})\) by virtue of the associativity and the fact that \(\mathrm{N}_{\mathbb{F}}(u)\) belongs to the centre of F.

An A-algebra isomorphic to a quaternion algebra is also called a quaternion algebra; if a basis of such an algebra has multiplication table (30), it is called a basis of type \((\alpha, \beta, \gamma)\) . By an abuse of language, a quaternion algebra is said to be of type \((\alpha, \beta, \gamma)\) when it has a basis of type \((\alpha, \beta, \gamma)\) .

When \(\beta = 0\) , formulae (30) and (31) simplify to

\[
\left\{ \begin{array}{l l} i ^ {2} = \alpha e, & \quad i j = k, \qquad i k = \alpha j, \\ j i = - k, & \quad j ^ {2} = \gamma e, \qquad j k = - \gamma i, \\ k i = - \alpha i, & \quad k j = \gamma i, \qquad k ^ {2} = - \alpha \gamma e, \end{array} \right.\tag{33}
\]

and

\[
\left\{ \begin{array} {c}\bar {u} = \rho e - \xi i - \eta j - \zeta k \\ \mathrm{T} _ {\mathrm{F}} (u) = 2 \rho \\ \mathrm{N} _ {\mathrm{F}} (u) = \rho^ {2} - \alpha \xi^ {2} - \gamma \eta^ {2} + \alpha \gamma \zeta^ {2}. \\ \end{array} \right.\tag{34}
\]

Then \((\alpha, \beta, \gamma)\) is replaced throughout by \((\alpha, \gamma)\) in the above expressions. It is immediate that the quaternion algebras of types \((\alpha, \gamma)\) and \((\gamma, \alpha)\) are isomorphic.

Note that formulae (32) show that F is not commutative when \(-1 \neq 1\) in A.

Taking A to be the field \(\mathbf{R}\) of real numbers and \(\alpha = \gamma = -1\) , \(\beta = 0\) , the corresponding algebra \(F\) is called the algebra of Hamiltonian quaternions and is denoted by \(\mathbf{H}\) . If \(u = \rho e + \xi i + \eta j + \zeta k\) ( \(\rho, \xi, \eta, \zeta\) in \(\mathbf{R}\) ) is an element \(\neq 0\) in \(\mathbf{H}\) , the formula \(u\bar{u} = \bar{u} u = \rho^2 + \xi^2 + \eta^2 + \zeta^2\) (cf.~(34)) shows that \(N(u) \neq 0\) in \(\mathbf{R}\) , so that \(u\) admits an inverse \(u^{-1} = N(u)^{-1}\bar{u}\) in \(\mathbf{H}\) and that \(\mathbf{H}\) is therefore a non-commutative field.

\begin{enumerate}
\def\labelenumi{(\arabic{enumi})}
\setcounter{enumi}{2}
\tightlist
\item
  If E is taken to be a quaternion algebra (cf.~Example 2), the Cayley extension of E defined by an element \(\delta \in A\) is in general non-associative (Proposition 5); it is called an octonion algebra over A (cf.~Appendix, no. 3).
\end{enumerate}

